\documentclass{article}
\usepackage{/home/user1/code/tex/sty}
\usepackage{amsfonts}
\usepackage{amsmath}
\usepackage{geometry}
\usepackage{graphicx}
\usepackage{hyperref}
\usepackage{enumitem}
\setlength{\parindent}{0pt}
\geometry{top=1in,bottom=1in,left=1in,right=1in}
\begin{document}
\begin{center}
\textbf{\Large{Woche 3 - Hausafgaben}}\\~\\
\begin{tabular}{l l}
Student&Agustin Romero\\
Email&ar1@stanford.edu\\
Instructor&Patric Di Dio De Marco\\
Course&GERLANG 3\\
Institution&Stanford University\\
Date&\today\\
\end{tabular}
\end{center}
\section*{Hausafgaben Woche 3 Tag 1}
\subsection*{181-5}
\begin{itemize}
\item Die Komparativ von ,,nah'' ist näher.
\item Die Komparativ von ,,warm'' ist wärmer.
\item Die Komparativ von ,,schnell'' ist scheller.
\item Die Komparativ von ,,ruhig'' ist ruhiger.
\item Die Komparativ von ,,billig'' ist billiger.
\item Die Komparativ von ,,kurz'' ist kürzer.
\end{itemize}
\begin{enumerate}[label=(\alph*)]
\item Der Flug nach New York kostet 800 €. 
\begin{itemize}
\item Das ist zu teuer. Wir müssen einen \underline{billigeren} Flug finden.
\end{itemize}
\item Der flug auf die Philippinen dauert 15 Stunden. Ist das zu lang?
\begin{itemize}
\item Ja, such wir doch ein \underline{kürzer} Urlaubsziel.
\end{itemize}
\item Mit unseren alten Fahrrädern dauert die Tour mindesetens drei Tage.
\ite{\item Dann leihen wir uns doch \underline{schneller} Räder.}
\item Willst du über Hamburg oder über Leipzig nach Berlin fahren?
\ite{\item Ich bin für die \underline{kürzer} Strecke.}
\item Ich möchte im Sommer nach Grönland reisen.
\ite{\item Ich würde lieber in ein \underline{warmer} Land fahren.}
\item Fahren wir doch eine Woche nach Berlin.
\ite{\item Ein \underline{näher} Ort wäre mir lieber.}
\end{enumerate}
\subsection*{181-6}
\begin{enumerate}[label={\alph*}]
\item Onkel Michael meint, früher gab es \underline{gringere Verkehrprobleme}.
\item Lilli meint, dass es heute \underline{preiswertere U-Bahnen gibt}.
\item Michael findet, dass es \underline{früher vernünftigere Freizeitbeschäftigungen}.
\item Lilli gibt ihrem Onkel den Rat, einen Computer zu kaufen, weil es keine \underline{einfachere Kommunikationsmöglichkeit gibt}.
\item Lilli findet, im Internet gibt es oft \underline{niedrigere Preise} als im Supermarkt.
\item Michael fragt, warum es \underline{zuverlässigere Verkehrsmittel} als grüher gibt, obwohl jeder einen Computer hat.
\end{enumerate}
\section*{Hausafgaben Woche 3 Tag 2}
\subsection*{182-2}
\itealph{\item Wir könnten \und{irgendwo} auf einem Parkplatz schlafen.
\ite{\item Ich will aber in einem guten Hotel schlafen.}
\item Ich möchte gern im Urlaub zu Hause bleiben. Müssen wir immer \und{irgendwohin} fahren?
\ite{\item Ja, ich will unbedingt nach Frankreich.}
\item Wir sollten \und{irgendetwas} als Geschenk mitbringen.
\ite{\item Stimmt. Vielleicht eine Flasche Rotwein?}
\item Wir können uns ja \und{irgendwann} nächste Woche treffen.
\ite{\item Gute Idee! Wir wär's mit Donnerstag?}
\item Ich glaube, wir sollten \und{irgendwen} nach dem Weg fragen.
\ite{\item Fragen wir doch den Mann dort drüben.}
\item \und{Irgendwie} müssen wir jetzt nach Frankfurt kommen.
\ite{\item Fahren wir doch per Autostopp.}}
\subsection*{183-5}
\itealph{\item Bevor wir ins Flugzeug steigen,
\ite{\item müssen wir das Gepäck einchecken.
\item müssen wir das Flugticket ausdrucken.
\item müssen wir das richtige Gate finden.}
\item Während wir fliegen, \ite{\item können wir die Alpen von oben sehen.
\item können wir nicht telefonieren.
\item bekommen wir vom Kapitän Informationen über das Wetter.}
\item Nachdem wir gelandet sind, \ite{\item müssen wir unser Gepäck abholen.
\item holt uns ein Mitarbeiter vom Hotel ab.}}
\section*{Hausafgaben Woche 3 Tag 3}
\subsection*{183-6}
\ite{\item Bevor ich innerhalb eine Straßentunnel fahren, solle ich prüfen ob Vehrkehrstau innerhalb stellt oder nicht.
\item Bevor ich gehen meinen Laboratorium hin, müsse ich schauen meinem Ausweis vor dem Eingang.
\item Bevor ich nach einem Zeitschrift eines neues Publikation senden, solle ich senden das nach allen meinen beziehenden Kollegen, um Felher zu prüfen.
\item Während ich fliegen in einem Flugzeug, ich immer schlafen, weil ich könne niemals lesen oder arbeiten.
\item Während ich prüfen Hausaufgaben von Studenten, müsse ich vergleichen sie mit die offizielen Antworten.
\item Nachdem ich gab eine Unterricht, solle ich waschen meinen Hände, denn ich habe auf mich viele Tafelkreide gestreut.
\item Nachdem ich habe eine Besprechung gefertigt, solle ich schreiben Besprechungsnotizen, deshalb wir können nicht so einfach vergessen was hat passiert.
}
\subsection*{183-8}
\itealph{\item Sie hatten bei dem Schloss getroffen. Er hat mit deiner Auto den Fluss entlang gefahren, dann habe er das Schloss gesehen. Er mag sich mehr anzusehen, dann in eine Stunde fangt seine Besichtigung an. Sie wohnte eine Tag in Hotel Lorenz. Sie an Mittagspause das Schloss passiert.

\item Zu dem Schloss sie hat dem Tor geht, und dann die Königstaße entlang, bis sie zu Parkstraße kam. Dort muss sie rechts gehen, ungefähr 50 meter weiter muss sie links in die Petersgasse abbiegen, und das Hotel auf der rechten Seite. Er kam aus seiner Auto, auf dem Besucherparkplatz vor der Staat. Er streckt Laubgasse entlang gegangen, und dann zuerst nach links an die Ecke gebogen, dabei er bei einem Café. Musste er wieder gehen zurück. Dort bin er bei dem Parkstraße gegangen, und links an dem Königstraße eingebogen.}
\subsection*{183-9}
\itealph{\item Vom Hotel Lorenz geht man zuerst nördlich bis Parkstraße, danach links auf Parkstraße für ein halb Häuserblock. Bevor Sie auf der Hauptplatz geht, leiten Sie rechts auf die Annawegstraße für knapp drei Häuserblock. Der Café Steiner ist gegenseitig der Königschloss.
\item Vom Parkplatz geht man osten auf dem Laubgasse, rechts an dem Annaweg, links auf dem Parkstraße, und zuletzt links auf dem Königstraße. Solle man nicht bei dem Hotel Lorenz passieren.
\item Zuerst solle man gehen ost auf dem Laubgasse. In kurze Zeit sehen Sie die Apotheke, an dem Ecke bei Laubgasse und Annaweg. Nachdem Sie seiner Medikament besorgt, sollen sie gehen südliche auf Annaweg. An dem Ecke des Annaweg und Parkstraße sollen Sie leiten rechts auf Parksträse, und vorbei die Hauptplatz an dem links finden Sie den Bank. Nachdem Sie seiner Geld getragt, sollen sie gehen zurück der Parkstraße, vorbei Annaweg, und links auf die Königstraße. Sehen Sie die Schloss an dem Ende des Königstraße. Nachdem Sie haben die Schloss sehen, sollen Sie gehen südliche auf Königstraße zurück, rechts auf Parkstraße, und rechts auf Annaweg, vorbei die Apotheke und Schloss, und sehen Sie die Café links dem Annaweg.}
\subsection*{184-4}
\itealph{\item Wenn meine Heimatstadt direkt an der Grenze zu Brasilien liegen würde, \und{würde ich jeden Abend Samba tanzen.}
\item Wenn es neben meinem Heimatdorf einen sehr hohen Berg geben würde, \und{würde ich jeden Tag Ski fahren.}
\item Wenn mein Haus direkt am Strand liegen würde, \und{würde ich ein großes Motorboot haben.}
}
\section*{Hausafgaben Woche 3 Tag 4}
Keine Hausaufgaben von dem Arbeitsbuch war gebietet.
\section*{Hausafgaben Woche 3 Tag 5}
Keine Hausaufgaben von dem Arbeitsbuch war gebietet.
\end{document}
